
This appendix documents the enterprise knowledge corpus extraction, deduplication, and preprocessing pipeline applied to the partner enterprise documentation corpus.

% ==============================================================================
% LIITE A: YRITYSKORPUKSEN (NORDEA) TIEDONKERUU- JA ESIKÄSITTELYPUTKI
% (Muistiinpanot suomeksi myöhempää englanninkielistä auki kirjoittamista varten)
% ==============================================================================
%
% 1. YLEISKUVA JA TAVOITE (Overview & Enterprise Context)
% ------------------------------------------------------------------------------
% - Tutkimuksessa tarvitaan aito ja realistinen suomenkielinen yritystason referenssiaineisto
%   (enterprise benchmark corpus), jotta RAG-optimointimenetelmien (M0–M4) tehokkuutta ja
%   laatua voidaan mitata tuotantovastaavassa ympäristössä ilman keinotekoisia synteettisiä datoja.
% - Kohteeksi valittiin Nordean julkinen henkilöasiakkaiden verkkopalvelu
%   (https://www.nordea.fi/henkiloasiakkaat/) ja siihen kiinteästi liittyvät PDF-ehdot/hinnastot.
% - Aineiston laajuus:
%     * 250 HTML-verkkosivua (henkilöasiakkaiden tilit, lainat, kortit, asuminen, säästäminen jne.)
%     * 50 PDF-dokumenttia (viralliset palveluehdot, hinnastot ja tuoteselosteet)
%     * Yhteensä 300 lähdedokumenttia.
% - Arkkitehtuurin toteutus:
%     * Integroitu osaksi Java 21 / Spring Boot 3 / Spring AI -testialustaa (com.dippa.ragbenchmark.ingest.nordea).
%     * Rinnalla itsenäinen Python-esikäsittelytyökalu (clean_nordea_markdown.py) komentorivikäyttöön.
%
% 2. TIEDONKERUUN KOMPONENTIT JA TEKNINEN TOTEUTUS (Crawling & Extraction)
% ------------------------------------------------------------------------------
% - NordeaCrawler:
%     * Leveyshakuun (Breadth-First Search, BFS) perustuva verkkoharavointi Jsoup-kirjastolla.
%     * Rajattu tiukasti vain "https://www.nordea.fi/henkiloasiakkaat/" -polun alle (ei eksy
%       muihin kieliversioihin kuten /sv/ tai /en/, eikä verkkopankin kirjautumissivuille).
%     * Siistii URL-parametrit (poistaa analytiikka- ja kampanjatagit kuten intcmp, icid, utm_*).
%     * Eettisyys ja kuormitussuojaus: 400 ms kohtelias viive (politeDelay) pyyntöjen välillä
%       sekä selkeä tutkimuskäyttöön yksilöity User-Agent -tunniste yhteystietoineen.
% - NordeaPdfExtractor:
%     * Lataa löydetyt PDF-linkit muistiin (kokorajoitus 20 MB ylisuurten tiedostojen varalta).
%     * Tekstinerotus suoritetaan Apache PDFBox -kirjastolla, ja teksti muotoillaan
%       rakenteelliseksi Markdowniksi otsikkotasoineen.
%
% 3. VIKASIETOISUUS JA UUDELLEENYRITYKSET (Resilience & Retry Mechanism)
% ------------------------------------------------------------------------------
% - Verkkoympäristössä esiintyy satunnaisia viiveitä, socket-aikakatkaisuja tai palvelinpään
%   kuormitushuippuja (HTTP 5xx). Alkuperäisessä ajossa 1–2 sivua saattoi ohittua häiriöiden vuoksi.
% - Toteutettiin eksponentiaaliseen viiveeseen (exponential backoff) perustuva uudelleenyritys:
%     * Enintään 3 yritystä (MAX_FETCH_RETRIES = 3).
%     * Viiveet yritysten välillä: 1000 ms, 2000 ms, 4000 ms.
%     * Mukana turvallinen säikeen keskeytyksen hallinta (!Thread.currentThread().isInterrupted()).
% - Mekanismi kattaa sekä HTML-sivujen noudon että PDF-tiedostojen HTTP-lataukset.
% - Vektorikantatallennuksessa (addInBatches) toteutettiin vastaava eräajon uudelleenyritys
%   mahdollisten PgVector-yhteyskatkosten varalta.
%
% 4. KOHINANPOISTO JA METODOLOGINEN PERUSTELU (Boilerplate Removal & Cleaning)
% ------------------------------------------------------------------------------
% - Tieteellinen perustelu (Data Preprocessing in RAG):
%     * Web-raakadatan jättäminen sellaisenaan ("in-the-wild") tuottaa strukturaalista kohinaa,
%       joka heikentää tutkimuskysymysten mukaista suorituskykyä:
%       - RQ1 & RQ2 (Latenssi ja kustannukset): Toistuvat evästebannerit ja linkkilistat syövät
%         LLM:n konteksti-ikkunaa, nostavat turhaan tokenikustannuksia ja pidentävät TTFT-latenssia.
%       - RQ3 (Laatu & Context Precision): Jos sama geneerinen 8-rivinen asiakaspalvelualatunniste
%         esiintyy jokaisella 245 sivulla, Small-to-Big -haussa syntyy satoja kilpailevia identtisiä
%         chunkkeja ("Kysy heti chatissa / Soita meille"), jotka syrjäyttävät varsinaisen oikean
%         yhteystietosivun hakutuloksista.
% - Deterministinen automaattinen siivous vs. manuaalinen käsin siivous:
%     * 300 tiedoston manuaalinen editointi käsin olisi tieteellisesti ongelmallista toistettavuuden
%       (reproducibility) kannalta ja lisäisi inhimillistä harhaa (selection bias / data snooping).
%     * Siivous toteutettiin täysin sääntöpohjaisena ja deterministisenä algoritmina, joka voidaan
%       dokumentoida ja toistaa koodipohjaisesti sekunneissa.
% - Puhdistetut kohinaelementit:
%     1. Evästemuuri- ja videovaroitukset:
%        "Hups! Jos haluat nähdä tällaista sisältöä, hyväksy ensin Nordean markkinointievästeet"
%        (Poistettu 85 esiintymää 29 tiedostosta; johtuvat iframe-video/podcast-upotuksista).
%     2. Saavutettavuus- ja ruudunlukijajäänteet:
%        "Avautuu uuteen ikkunaan" (Poistettu 26 esiintymää 17 tiedostosta otsikoista ja linkeistä).
%     3. Toistuva asiakaspalvelualatunniste (Standard Help Footer Block):
%        "## Apua asiointiin \n ### Henkilöasiakkaat \n - Kysy heti chatissa \n - Usein kysytyt \n - Soita meille \n - Nordea Omaposti"
%        (Poistettu 245 tiedostosta 250:stä. Tärkeää: varsinaisen yhteystietosivun oma asiasisältö,
%        aukioloajat ja puhelinnumerot säilytettiin täysin koskemattomina).
%     4. Rakenteellinen normalisointi:
%        Orpojen tyhjien otsikoiden poisto ja moninkertaisten tyhjien rivien tiivistäminen (max \n\n).
%     5. Metadata / YAML Frontmatter:
%        Tiedostojen alussa oleva metadata (url, title, type, crawled_at) säilytettiin 100 % ehjänä.
%
% 5. TIETOPUTKEN VAIHEISTUS JA OHJAUS (Pipeline Orchestration)
% ------------------------------------------------------------------------------
% - REST-rajapinta: POST /api/ingest/nordea
% - Parametrit (NordeaIngestRequest):
%     * action: "all" (täysi putki), "crawl" (vain verkkoharavointi levylle),
%       "clean" (vain levytiedostojen siivous), "ingest" (vain levytiedostojen indeksointi kantaan),
%       "s2b-baseline" (suora kannan sisäinen muunnos baselinesta S2B-tauluihin).
%     * mode: "s2b" (small-to-big: vector_store_nordea_large & _small) tai
%       "original" (baseline: vector_store_nordea).
%     * clearDatabase: true/false (truncoi kohdevektoritaulut ennen ajoa).
%     * clearFiles: true/false (poistaa vanhat .md-tiedostot levyltä ennen uutta crawlia).
% - Erottelu offline- ja online-vaiheiden välillä:
%     * Kun aineisto on kerran haravoitu levylle, se voidaan siivota ja indeksoida uudelleen
%       kantaan sekunneissa ilman hidasta tai kuormittavaa verkkocrawlausta (action: "ingest").
%
% 6. VEKTORITAULUT JA CHUNKKAAUSSTRATEGIAT (Vector Storage & Splitting)
% ------------------------------------------------------------------------------
% - M0 Baseline:
%     * Taulu: vector_store_nordea
%     * Chunk-koko: 512 tokenia (TokenTextSplitter).
% - M1 Small-to-Big (S2B):
%     * Hakuindeksi (Child): vector_store_nordea_small, chunk-koko 128 tokenia.
%       Takaa tarkan semanttisen haun ja korkean Context Precisionin.
%     * Kontekstilaajennus (Parent): vector_store_nordea_large, chunk-koko 1024 tokenia.
%       Palauttaa generointivaiheeseen eheän kontekstin ja ratkaisee semanttisen pirstaloitumisen.
%     * Kytkentä: Jokainen child-chunk sisältää metatiedoissaan viitteen parent_chunk_id.
% - Baseline -> S2B -migraatio (s2b-baseline):
%     * Mahdollistaa Small-to-Big -indeksien rakentamisen suoraan olemassa olevan
%       vector_store_nordea -kannan pohjalta ilman levy- tai verkkotoimenpiteitä.
% 7. VERTAILUASETELMA: OPEN-CORPUS VS. CLOSED-CORPUS (Enterprise Closed & MuSiQue Parity)
% ------------------------------------------------------------------------------
% - Tausta ja metodologinen ongelma (The Confounding Variable Problem):
%     * MuSiQue-vertailuaineisto ingestoitiin alun perin suljettuna korpuksena (closed corpus):
%       vektorikantaan ladattiin vain 100 benchmark-kysymyksen mukana tulevat 20 kandidaattikappaletta
%       per kysymys, jolloin deduplikoinnin jälkeen kannassa oli vain ~1 700 kappaletta.
%     * Nordea-yritysaineiston osalta oli sen sijaan crawlattu ja indeksoitu koko verkkopalvelu
%       (open corpus: 250 sivua + 50 PDF:ää, tuhansia chunkkeja).
%     * Tämä muodosti tutkimukseen vakavan sekoittavan muuttujan (confounding variable):
%       - Suora vertailu MuSiQuen ja Nordean välillä vääristyi, sillä hakuavaruuden vaikeusaste oli
%         täysin eri ("neula heinäsuovassa" vs. kontrolloitu 20 kappaleen häirikköjoukko).
%       - Esimerkiksi Context Precision ja Context Recall olivat Nordeassa matalampia puhtaasti
%         laajemman hakuavaruuden vuoksi, ei järjestelmän heikomman suorituskyvyn takia.
%       - Optimointimoduulien hyödyt vääristyivät: Re-ranking (M2) tuottaa suhteettoman suuren
%         lisäarvon laajan korpuksen kohinassa verrattuna 1 700 kappaleen rajattuun aineistoon.
% - Ratkaisu: Kolmitahoinen koesarja (Three-way Experimental Paradigm / Vaihtoehto C):
%     * Toteutettiin Nordealle vastaava suljettu korpus (enterprise_closed), jotta aineistojen
%       välille saadaan tieteellisesti 1:1 vertailukelpoinen testiasetelma.
%     * Samalla säilytettiin täysi avoimen korpuksen aineisto ja sen ajotulokset koskemattomina.
%     * Diplomityöhön muodostuu poikkeuksellisen vahva vertaileva tutkimusasetelma:
%       1. MuSiQue (EN Closed) <-> Nordea Closed (FI Closed):
%          -> Puhdas kieli- ja toimialariippuvuuden vertailu ilman hakuavaruuden kokoeroa.
%       2. Nordea Closed (FI Closed) <-> Nordea Open (FI Open):
%          -> Korpuksen laajuuden ja kohinan vaikutus RAG-optimointiin (Pareto-rintaman
%             muutos ja M2 Re-rankingin kriittisyys reaalimaailman tuotantoskaalassa).
% - Tekninen toteutus (Technical Implementation):
%     * Lähdeaineisto: nordea_100_final.jsonl (sisältää 100 kysymystä x 20 kappaletta = 2 000 kpl).
%     * Deduplikointi: Identtisellä MuSiQue-logiikalla (title + "\n" + paragraph_text) saatiin
%       tarkalleen 1 333 uniikkia kappaletta (hyvin lähellä MuSiQuen ~1 700 kappaleen mittakaavaa).
%     * Uudet erilliset pgvector-taulut (ei kosketa olemassa oleviin):
%       1. vector_store_nordea_closed:
%          Perushaku (C0 Baseline, C2 Reranker, C3 Prompt Cache, C4 Router) 1 333 kappaleelle.
%       2. vector_store_nordea_closed_large:
%          Small-to-Big (M1) kontekstilaajennus (parent-chunkit 1024 tokenia).
%       3. vector_store_nordea_closed_small:
%          Small-to-Big (M1) tiheä hakuindeksi (child-chunkit 128 tokenia, parent_id -linkitys).
%     * Dynaaminen kyselyreititys:
%       RagQueryService reitittää kyselyn taulun nimen mukaan:
%       Jos tableName = "vector_store_nordea_closed", M1-kyselyt ohjataan automaattisesti
%       pariin vector_store_nordea_closed_small -> vector_store_nordea_closed_large.
%     * Eristys ja hakemistorakenne:
%       Kaikki suljetun korpuksen ajot, taulukot ja kuvaajat eriytetään omaan alikansioonsa:
%       - Tulokset: rag-evaluator/runs/enterprise_closed/
%       - Taulukot: thesis/tables/enterprise_closed/
%       - Kuvaajat: thesis/figures/enterprise_closed/
%       Tämä takaa, että alkuperäiset open-corpus -tulokset (enterprise/) eivät ylikirjoitu.
% ==============================================================================
